% Part: computability
% Chapter: machines-computations
% Section: halting-problem

\documentclass[../../../include/open-logic-section]{subfiles}

\begin{document}

\olfileid{tur}{und}{hal}
\olsection{थांबण्याची समस्या}

\begin{explain}
ट्यूरिंग यंत्रांच्या वर्णनांचे एक प्रगणन आपण निश्चित
केले आहे असे समजा. त्यामुळे प्रत्येक ट्यूरिंग
यंत्राला एक \emph{निर्देशांक} मिळतो: ट्यूरिंग
यंत्रांच्या वर्णनांच्या $M_1$, $M_2$, $M_3$,
\dots{} या प्रगणनातील त्याचे स्थान.

ट्यूरिंग-संगणनक्षम नसलेली फलने असलीच पाहिजेत,
हे आपल्याला माहीत आहे: ट्यूरिंग यंत्रांच्या
वर्णनांचा---आणि म्हणून यंत्रांचा---संच
!!{enumerable} म्हणजे गणनीय आहे; पण
$\Nat$ पासून $\Nat$ मध्ये जाणाऱ्या सर्व
फलनांचा संच तसा नाही. मात्र संगणनक्षम
नसलेल्या विशिष्ट फलनांची उदाहरणेही देता
येतात. त्यांपैकी एक म्हणजे थांबण्याचे फलन.
\end{explain}

\begin{defn}[थांबण्याचे फलन]
 \emph{थांबण्याचे फलन}~$h$ असे परिभाषित करतात:
\[
h(e,n) =
\begin{cases}
  \text{0} & \text{जर यंत्र~$M_e$ आदान $n$ वर थांबत नसेल तर} \\
  \text{1} & \text{जर यंत्र~$M_e$ आदान $n$ वर थांबत असेल तर}
\end{cases}
\]
\end{defn}

\begin{defn}[थांबण्याची समस्या]
\emph{थांबण्याची समस्या} म्हणजे (कोणत्याही $e$, $n$
साठी) ट्यूरिंग यंत्र~$M_e$ हे $n$ रेघांच्या
आदानावर थांबते का हे ठरवण्याची समस्या.
\end{defn}

\begin{explain}
~$h$ हे ट्यूरिंग-संगणनक्षम नाही, हे आपण
त्याच्याशी संबंधित~$s$ हे फलन ट्यूरिंग-संगणनक्षम
नाही असे दाखवून सिद्ध करू. हा पुरावा कोणत्याही
ट्यूरिंग यंत्राचे संगणन शिस्तबद्ध ट्यूरिंग
यंत्राने करता येते (\olref[mac][dis]{sec}),
आणि दोन ट्यूरिंग यंत्रे जोडून एक यंत्र
बनवता येते (\olref[mac][cmb]{sec}),
या दोन गोष्टींवर आधारलेला आहे.
\end{explain}

\begin{defn} फलन~$s$ पुढीलप्रमाणे परिभाषित आहे:
\[
s(e) =
\begin{cases}
  \text{0} & \text{जर यंत्र~$M_e$ आदान $e$ वर थांबत नसेल तर} \\
  \text{1} & \text{जर यंत्र~$M_e$ आदान $e$ वर थांबत असेल तर}
\end{cases}
\]
\end{defn}

\begin{lem}
फलन~$s$ ट्यूरिंग-संगणनक्षम नाही.
\end{lem}

\begin{proof}
विरोधासाठी फलन~$s$ ट्यूरिंग-संगणनक्षम आहे असे
समजा. मग~$s$ संगणित करणारे ट्यूरिंग यंत्र~$S$
असेल. व्यापकतेला बाधा न आणता, $S$ शिस्तबद्ध
आहे असे मानता येते. विशेषतः, ते थांबते
तेव्हा पहिल्या घराकडे पाहत असते.
%READERNOTE{OLTUR-008 — गोठवलेल्या इंग्रजी पुराव्यात पहिले घर पाहत थांबणे म्हणजेच शिस्तबद्ध असणे असे कंसात म्हटले आहे. आधीच्या व्याख्येनुसार शिस्तबद्ध यंत्राला याशिवाय खास थांबण्याची अवस्था, टोकाचे चिन्ह राखणे आणि घर ० वरून डावीकडे न जाणे या अटीही आहेत. येथे आधीच्या रूपांतरानुसार पूर्ण शिस्तबद्ध यंत्र गृहीत धरले आहे; त्यातील पहिल्या घरावर थांबण्याचा गुणधर्म या जोडणीसाठी वापरला जातो.}
या यंत्राला दुसरे यंत्र~$J$ जोडता येते.
त्याला आदान~$0$ दिल्यास ते थांबते
(म्हणजे आरंभीच्या अवस्थेत, फितीच्या टोकाच्या
चिन्हाच्या उजवीकडील घराकडे पाहताना त्याला
$\TMblank$ दिसल्यास); अन्यथा ते उजवीकडे
चालतच राहते आणि कधीच थांबत नाही.
$S$ ला~$J$ जोडून तयार केलेले
$S \concat J$ हे ट्यूरिंग यंत्र आहे;
म्हणून ते कोणत्यातरी~$e$ साठी $M_e$
आहे (म्हणजे प्रगणनात कुठेतरी येते).
$M_e$ हे $e$ $\TMstroke$ चिन्हांच्या
आदानावर सुरू करा. दोनच शक्यता आहेत:
$M_e$ थांबेल किंवा थांबणार नाही.
\begin{enumerate}
\item $M_e$ हे $e$ $\TMstroke$ चिन्हांच्या
  आदानावर थांबते असे समजा. मग $s(e)
  = 1$. म्हणून~$e$ वर सुरू केलेले $S$
  थांबते आणि फितीवर फलित म्हणून एकच
  $\TMstroke$ ठेवते. मग $J$ फितीवरील
  $\TMstroke$ चिन्हापासून सुरू होते. त्या
  बाबतीत $J$ थांबत नाही. पण $M_e$ हे
  $S \concat J$ यंत्रच आहे; म्हणून $S$
  नंतर~$J$ जे करेल तेच त्याने करायला
  हवे (म्हणजे या बाबतीत उजवीकडे चालत
  राहून कधीच न थांबणे). म्हणून $M_e$
  हे $e$ $\TMstroke$ चिन्हांच्या आदानावर
  थांबू शकत नाही.

\item आता $M_e$ हे $e$ $\TMstroke$
  चिन्हांच्या आदानावर थांबत नाही असे
  समजा. मग $s(e) = 0$ आणि आदान~$e$
  वर सुरू केलेले $S$ रिकामी फीत ठेवून
  थांबते. रिकाम्या फितीवर सुरू केलेले
  $J$ लगेच थांबते. पुन्हा, $M_e$ हे
  $S$ नंतर~$J$ जे करेल तेच करते;
  म्हणून $M_e$ हे $e$ $\TMstroke$
  चिन्हांच्या आदानावर थांबलेच पाहिजे.
\end{enumerate}
दोन्ही बाबतीत आपल्या गृहीतकाला विरोध
मिळतो. त्यामुळे असे ट्यूरिंग यंत्र~$S$
असू शकत नाही: $s$ ट्यूरिंग-संगणनक्षम नाही.
\end{proof}

\begin{thm}[थांबण्याच्या समस्येची असमाधेयता]
\ollabel{thm:halting-problem} थांबण्याची समस्या
सोडवता येत नाही; म्हणजे फलन~$h$
ट्यूरिंग-संगणनक्षम नाही.
\end{thm}

\begin{proof}
$h$ हे ट्यूरिंग-संगणनक्षम आहे, आणि समजा
ट्यूरिंग यंत्र~$H$ ते संगणित करते.
मग~$H$ वापरून~$s$ संगणित करणारे ट्यूरिंग
यंत्र बनवता येईल: प्रथम आदानाची एक प्रत
बनवा (दोन्ही प्रतींमध्ये~$\TMblank$ चिन्ह
ठेवा). मग पुन्हा आरंभी जा आणि~$H$
चालवा. पहिले काम करणारे यंत्र स्पष्टपणे
बनवता येते (\cref{tur:mac:dis:prob:copier}
पाहा); आणि $H$ अस्तित्वात असेल, तर
ते अशा प्रतिकृती तयार करणाऱ्या यंत्राला
जोडून $M_e$ हे आदान~$e$ वर थांबते
का हे ठरवणारे, म्हणजे~$s$ संगणित
करणारे, नवे यंत्र मिळेल. पण असे यंत्र
असू शकत नाही, हे आपण आधीच दाखवले
आहे. म्हणून~$h$ देखील ट्यूरिंग-संगणनक्षम
नाही.
\end{proof}

\begin{prob}
तीन-रेघांची थांबण्याची (3-Halt) समस्या
म्हणजे अन्यथा रिकाम्या असलेल्या फितीवर
तीन $\TMstroke$ चिन्हांचे आदान दिल्यावर
स्वैरपणे निवडलेले ट्यूरिंग यंत्र थांबते
की नाही हे ठरवणारी निर्णय प्रक्रिया
देण्याची समस्या. 3-Halt समस्या
सोडवता येत नाही हे सिद्ध करा.
\end{prob}

\begin{prob}
ट्यूरिंग यंत्र आणि आदान यांच्या
$M_e$ आणि~$n$ या जोड्यांसाठी,
$e \neq n$ असताना थांबण्याची समस्या
सोडवता येत असेल, तर $e = n$
असलेल्या बाबतीतही ती सोडवता
येते, हे दाखवा.
\end{prob}

\begin{prob}
आदान~$e$ ही संख्या असेल आणि ती
सर्व ट्यूरिंग यंत्रांच्या प्रगणनाद्वारे
यंत्र~$M_e$ ओळखून देत असेल,
तर थांबण्याची समस्या सोडवता
येत नाही, हे आपण सिद्ध केले.
त्याऐवजी \olref[tur][und][enu]{sec}
मधील ट्यूरिंग यंत्रांची वर्णनेच
थेट आदान म्हणून दिली तर काय?
निर्देशांकांऐवजी ट्यूरिंग यंत्रांची
वर्णने आदान म्हणून घेऊन थांबण्याच्या
समस्येचा निर्णय करणारे ट्यूरिंग
यंत्र असू शकते का? का किंवा
का नाही, ते स्पष्ट करा.
\end{prob}

\begin{prob} पुढीलप्रमाणे परिभाषित
  \emph{आंशिक} फलन~$s'$ हे
  ट्यूरिंग-संगणनक्षम \emph{आहे},
  हे दाखवा:
  \[
  s'(e) =
  \begin{cases}
    \text{1} & \text{जर यंत्र~$M_e$ आदान $e$ वर थांबत असेल तर}\\
    \text{अपरिभाषित} & \text{जर यंत्र~$M_e$ आदान $e$ वर थांबत नसेल तर}
  \end{cases}
  \]
\end{prob}

\end{document}
